\chapter{Gładkie przeszkody Donaldsona w wymiarze cztery}
\label{ch:donaldson}
\index{twierdzenie Donaldsona}
\index{równanie antysamodualności}

Klasyfikacja Freedmana z Rozdziału~\ref{ch:freedman-klasyfikacja}
dotyczy kategorii topologicznej. Pytanie, które z jej przykładów
można wygładzić, wymaga innego narzędzia. Donaldson wykorzystał
połączenia na wiązkach nad gładką rozmaitością czterowymiarową
i analizę ich krzywizny. Poniżej wyprowadzamy równanie
antysamodualności oraz jego prosty rachunek energii. Głębokie
twierdzenie o przestrzeni rozwiązań i wynikającą z niego
diagonalizację podajemy jako wynik zewnętrzny.

\section{Krzywizna połączenia i gwiazda Hodge'a}
\label{sec:donaldson-asd}

Niech $X$ będzie zamkniętą, zorientowaną, gładką
rozmaitością Riemanna wymiaru $4$, a $P\to X$ główną
wiązką o grupie $\mathrm{SU}(2)$. Grupę i jej algebrę
opisuje Dodatek~\ref{sec:app-grupa-u}; połączenia
omawialiśmy w Rozdziale~\ref{ch:koneksje}, a gwiazdę Hodge'a
w Sekcji~\ref{sec:hodge-four-manifolds}.
W lokalnej trywializacji połączenie opisuje
$1$-forma macierzowa $A$ o wartościach w
$\mathfrak{su}(2)$, a jego krzywiznę definiuje wzór
\begin{equation}\label{eq:donaldson-curvature}
 F_A=dA+A\wedge A.
\end{equation}
Iloczyn w drugim składniku łączy iloczyn zewnętrzny form
z mnożeniem macierzy. Dla pól stycznych $u,v$ mamy więc
$(A\wedge A)(u,v)=[A(u),A(v)]$.
Zmiana trywializacji przez $g:X\to\mathrm{SU}(2)$ daje
$A^g=g^{-1}Ag+g^{-1}dg$ oraz
$F_{A^g}=g^{-1}F_Ag$. Ostatnia równość wynika po
rozwinięciu $d(g^{-1})=-g^{-1}(dg)g^{-1}$:
wyrazy z $dg$ znoszą się parami. Dlatego warunek
na krzywiznę nie zależy od wybranej trywializacji.

W wymiarze cztery gwiazda Hodge'a spełnia
$*^2=1$ na $2$-formach. Dzielimy więc krzywiznę na
\[
 F_A^+=\tfrac12(F_A+*F_A),\qquad
 F_A^-=\tfrac12(F_A-*F_A).
\]

\begin{definicja}[Połączenie antysamodualne]
\label{def:donaldson-asd}
Połączenie $A$ na $P$ jest \emph{antysamodualne},
gdy $F_A^+=0$, czyli $*F_A=-F_A$.
Przestrzeń takich połączeń modulo działanie
gładkich przekształceń cechowania nazywamy
\emph{przestrzenią moduli ASD}.
\end{definicja}

Intuicja jest następująca: gwiazda rozdziela
krzywiznę na dwie niezależne połowy, a równanie
zeruje jedną z nich. Płaskie połączenie
$F_A=0$ jest prostym przykładem rozwiązania.
Nie twierdzimy, że na dowolnej wiązce istnieje
rozwiązanie niepłaskie; zależy to od topologii
wiązki i metryki.

\section{Energia i liczba Cherna}
\label{sec:donaldson-energia}

Na $\mathfrak{su}(2)$ przyjmujemy dodatni iloczyn
$\langle B,C\rangle=-\operatorname{tr}(BC)$.
Niech $\|F_A\|_{L^2}^2=\int_X|F_A|^2\,d\mathrm{vol}$.
Przy konwencji użytej niżej druga liczba Cherna wynosi
\begin{equation}\label{eq:donaldson-c2}
 k=c_2(P)[X]
   =\frac{1}{8\pi^2}\int_X\operatorname{tr}(F_A\wedge F_A)
   \in\Z.
\end{equation}
Całkowitość i niezależność od połączenia są tu
zastosowaniem twierdzenia Cherna--Weila z
Rozdziału~\ref{ch:chern-weil}; sam rachunek energii
możemy przeprowadzić bez analizy przestrzeni moduli.

\begin{stwierdzenie}[Rachunek energii ASD]
\label{prop:donaldson-energy}
Dla powyższych konwencji
\[
 \|F_A\|_{L^2}^2
 =8\pi^2k+2\|F_A^+\|_{L^2}^2.
\]
W szczególności dla połączenia antysamodualnego
energia wynosi $8\pi^2k$ i $k\geq0$.
\end{stwierdzenie}
\begin{proof}
Ponieważ $F_A^+$ i $F_A^-$ są ortogonalne,
\[
 |F_A|^2=|F_A^+|^2+|F_A^-|^2.
\]
W zorientowanej bazie ortonormalnej dla skalarnej
$2$-formy $\omega$ zachodzi
$\omega\wedge\omega=\langle\omega,*\omega\rangle
\,d\mathrm{vol}$. Stosując tę równość składowo
do macierzy i znak z definicji iloczynu na
$\mathfrak{su}(2)$, otrzymujemy
\[
 \operatorname{tr}(F_A\wedge F_A)
 =\bigl(|F_A^-|^2-|F_A^+|^2\bigr)
   \,d\mathrm{vol}.
\]
Równość~\eqref{eq:donaldson-c2} daje więc
$8\pi^2k=\|F_A^-\|_{L^2}^2-\|F_A^+\|_{L^2}^2$.
Dodanie $2\|F_A^+\|_{L^2}^2$ kończy dowód.
\end{proof}

\section{Twierdzenie o diagonalizacji i jego zakres}
\label{sec:donaldson-diagonalizacja}

Przestrzeń moduli ASD nie jest zwykłą rozmaitością
zadaną skończoną liczbą równań. Trzeba wybrać
warunek cechowania, uzyskać lokalny model eliptyczny,
zbadać rozwiązania redukowalne oraz utratę zwartości
przez koncentrację krzywizny. Zwartość Uhlenbeck,
usuwanie osobliwości i analiza końców przestrzeni moduli
są odrębnymi głębokimi wynikami. Rozdział nie zawiera
ich dowodów.

\begin{twierdzenie}[Donaldson: diagonalizacja formy określonej;
wynik zewnętrzny]
\label{thm:donaldson-diagonalizacja}
Niech $X$ będzie zamkniętą, zorientowaną, gładką,
po prostu spójną rozmaitością wymiaru $4$.
Jeśli całkowita forma przecięcia $Q_X$ na
$H_2(X;\Z)/\operatorname{Tors}$ jest dodatnio określona,
to jest całkowicie izometryczna z
$\langle1\rangle^{\oplus b_2(X)}$.
Jeśli jest ujemnie określona, jest izometryczna z
$\langle-1\rangle^{\oplus b_2(X)}$.
\end{twierdzenie}

Pierwszą wersję otrzymuje się z drugiej po zmianie
orientacji. Źródłem jest
\href{https://doi.org/10.4310/jdg/1214437665}
{S.~K.~Donaldson, \emph{An application of gauge theory to
four-dimensional topology}, J.~Differential Geom. 18 (1983),
twierdzenie~A}. Dla zrozumienia metody: po odwróceniu
orientacji można badać wiązkę z $c_2=1$ nad rozmaitością
o ujemnie określonej formie. W odpowiedniej przestrzeni
moduli pojawiają się rozwiązania redukowalne i koniec
odpowiadający koncentracji krzywizny. Ich zgodność
topologiczna narzuca wektory kwadratu $-1$ tworzące
bazę formy. Uzasadnienie tego ostatniego zdania wymaga
właśnie analizy wymienionej powyżej; jest opisem struktury
dowodu Donaldsona, a nie dowodem twierdzenia.

\begin{wniosek}[Forma $E_8$ nie pochodzi od gładkiej
zamkniętej $4$-rozmaitości]
\label{cor:donaldson-e8}
Nie istnieje zamknięta, gładka, po prostu spójna
rozmaitość $4$-wymiarowa o formie przecięcia $E_8$.
Topologiczna rozmaitość Freedmana o tej formie
nie dopuszcza struktury gładkiej.
\end{wniosek}
\begin{proof}
Z obliczenia w
Stwierdzeniu~\ref{prop:freedman-E8-properties} forma
$E_8$ jest dodatnio określona, parzysta i ma rangę $8$.
Twierdzenie~\ref{thm:donaldson-diagonalizacja}
wymuszałoby izometrię z $I_8$. Ta ostatnia forma
jest nieparzysta, bo wektor bazowy ma kwadrat $1$;
parzystość zachowują izometrie całkowite. To sprzeczność.
Istnienie topologicznej rozmaitości o tej formie jest
częścią zewnętrznego twierdzenia Freedmana przyjętego
w Rozdziale~\ref{ch:freedman-klasyfikacja}.
\end{proof}

Ten wniosek pokazuje precyzyjną granicę między
klasyfikacją topologiczną a gładką. Nie wynika z niego
jeszcze istnienie egzotycznej $\R^4$: potrzebna jest
osobna konstrukcja otwartej rozmaitości i zwarta
przeszkoda rozróżniająca struktury. Przeprowadzimy
wynikającą z niej indukcję w następnym rozdziale.

\section*{Dalsza lektura}
\href{https://doi.org/10.4310/jdg/1214437665}
{S.~K.~Donaldson, \emph{An application of gauge theory to
four-dimensional topology}, J.~Differential Geom. 18 (1983)},
zawiera twierdzenie o diagonalizacji.
\href{https://doi.org/10.1090/S0273-0979-1983-15090-5}
{S.~K.~Donaldson, \emph{Self-dual connections and the topology
of smooth 4-manifolds}, Bull.~Amer.~Math.~Soc. 8 (1983)},
przedstawia metodę cechowania i kontekst wyniku.
